\section{Fundamentação Teórica}
\label{sec:fundamentacao}

\subsection{Conjuntos disjuntos e operações}

Seja um universo $\{0,1,\ldots,n-1\}$ particionado em conjuntos disjuntos. Cada conjunto é representado por um elemento identificador (raiz). A operação $\texttt{Find}(x)$ devolve o representante do conjunto que contém $x$; $\texttt{Union}(a,b)$ funde os conjuntos que contêm $a$ e $b$, caso sejam distintos. Duas consultas $\texttt{Find}(x)$ e $\texttt{Find}(y)$ retornam o mesmo valor se, e somente se, $x$ e $y$ pertencem ao mesmo conjunto.

\subsection{Floresta com união ingênua}

Uma forma clássica de representar conjuntos disjuntos consiste em manter uma floresta por meio de um vetor $\texttt{parent}$, no qual cada elemento aponta para seu pai e toda raiz satisfaz $\texttt{parent}[r]=r$. A operação $\texttt{Find}$ percorre esses ponteiros até alcançar a raiz, enquanto $\texttt{Union}$ localiza as raízes de $a$ e $b$ e faz uma delas apontar para a outra. Sem qualquer critério de balanceamento, a altura da árvore pode crescer linearmente, levando $\texttt{Find}$ a custo $\mathcal{O}(n)$ no pior caso e, portanto, sequências de $m$ operações podem atingir $\mathcal{O}(mn)$.

\subsection{União por rank}

Na heurística de união por rank, associa-se a cada nó um \emph{rank}, que funciona como majorante da altura da subárvore enraizada naquele ponto. Em $\texttt{Union}$, a raiz de menor rank é anexada à de maior rank; em caso de empate, incrementa-se o rank da nova raiz. Com essa estratégia, a altura das árvores permanece em $\mathcal{O}(\log n)$, de modo que cada $\texttt{Find}$ ou $\texttt{Union}$ custa $\mathcal{O}(\log n)$ no pior caso.

\subsection{Compressão de caminho e resultado de Tarjan}

Na compressão de caminho, cada nó visitado durante $\texttt{Find}$ passa a apontar diretamente para a raiz do conjunto, o que achata a estrutura ao longo do tempo. Quando essa heurística é combinada à união por rank, a análise amortizada sobre uma sequência de $m$ operações (com $m \geq n$) mostra que o tempo total é $\mathcal{O}(m\,\alpha(n))$, onde $\alpha(n)$ é a inversa de Ackermann, função que cresce tão lentamente que, para todos os valores de $n$ encontrados em aplicações práticas, $\alpha(n)\leq 4$ \cite{tarjan1975efficiency}.

\subsection{Interpretação prática de $\alpha(n)$}

A função de Ackermann cresce muito rapidamente, enquanto sua inversa $\alpha(n)$ cresce de forma extremamente lenta. Por isso, embora a complexidade $\mathcal{O}(\alpha(n))$ não seja constante em sentido estrito, na prática as operações se comportam quase como tempo constante amortizado para tamanhos de instância usuais. Esse comportamento justifica a presença recorrente do DSU em algoritmos de grafos, como Kruskal.
